\documentclass[10pt,a4paper]{amsart}
\usepackage[margin=19mm]{geometry}
\usepackage{amsmath,amssymb,mathtools}
\usepackage[scr=rsfs,cal=euler]{mathalfa}
\usepackage{xfrac}
\usepackage[unicode,hidelinks,bookmarks=false]{hyperref}
\newcommand{\ie}{i.e.\ }
\newcommand{\cf}{cf.\ }
\newcommand{\resp}{respectively\ }
\newcommand{\Subsection}{\subsection}
\providecommand{\nobreakdash}{\nobreak\hbox{-}}
\setlength{\emergencystretch}{2em}
\multlinegap0pt
\usepackage{bm}
\usepackage{bbm}
\usepackage{dsfont}
\usepackage{xspace}
\usepackage{cancel}
\usepackage{mathtools}
\usepackage{amsthm}
\makeatletter
\@ifundefined{theorem}{\newtheorem{theorem}{Theorem}}{}
\@ifundefined{lemma}{\newtheorem{lemma}[theorem]{Lemma}}{}
\makeatother
\makeatletter
\newcommand{\N}{\mathbb{N}}
\newcommand{\R}{\mathbb{R}}
\newcommand{\Z}{\mathbb{Z}}
\renewcommand{\d}{\,\textnormal{d}}
\expandafter\ifx\csname keywords\endcsname\relax
\newenvironment{keywords}{\begin{center}Keywords:}{\end{center}}
\fi
\DeclareMathOperator{\vol}{vol}
\makeatother
\title[A note on the complexity of two-stage stochastic linear opti]{A note on the complexity of two-stage stochastic linear optimization with small second stage}
\author{Christoph Buchheim}
\date{Source: arXiv:2605.25028v1 (CC BY 4.0)}
\begin{document}
\maketitle

\noindent\emph{Source and licence.} A note on the complexity of two-stage stochastic linear optimization with small second stage. Version arXiv:2605.25028v1 (CC BY 4.0), distributed under the Creative Commons Attribution 4.0 International licence, as recorded in the arXiv licence metadata for this exact version and shown on the abstract page https://arxiv.org/abs/2605.25028v1. Full rights notice: licenses/arxiv-2605.25028-CC-BY-4.0.txt.\par\smallskip

\noindent\textit{Editorial scope.} Counted unit: the Lemma whose source label is \texttt{lem:pseudo\_vol} in the article (source0.tex lines 804--807, source label \texttt{lem:pseudo\_vol}), delivered together with the article's own complete original proof of it (source0.tex lines 808--846, one \texttt{proof} environment of 39 lines). No mathematics is written, repaired or re-derived: statement and proof are the article's own text line for line. Required context retained uncounted: lem:p\_poly (lines 777--781). Every definition, display and label of those lines is retained unchanged. Omitted from this package and not redistributed: 1--776, 782--803, 847--866. Nothing inside a retained excerpt is omitted or altered: every retained body is delivered exactly as the article's lines are written, so no omission receipt of any kind accompanies this package. Citations: the retained excerpts carry no citation command at all, so no bibliography entry of the article is transcribed and no reference is redistributed. Reference adaptation: none was needed and none was made; every reference inside the delivered unit resolves inside it, so no unresolved reference or citation is produced. Printed slips kept literal: no printed word, symbol, hypothesis or number of the counted statement or of its proof was altered. Layout and class: the article's own class is \texttt{article} with packages including fullpage, amsthm, amsmath, amssymb, amsfonts, tikz; the wrapper is the lane's pinned \texttt{amsart} editorial wrapper at 10pt on A4, which restates the theorem-like environments the retained text needs and adds the article's own macro definitions that the retained text needs (\texttt{\textbackslash N}, \texttt{\textbackslash R}, \texttt{\textbackslash Z}, \texttt{\textbackslash d}, \texttt{\textbackslash vol}), with extra wrapper packages (bm, bbm, dsfont, xspace, cancel, mathtools). The delivered numbering is the unit's own. Assets: no retained span contains a figure environment, a graphics-inclusion command, a PDF-page inclusion, a table or a TikZ picture; no image file is copied into this package, the package declares no asset dependency and no image file is redistributed with it. Licence: version arXiv:2605.25028v1 of this article is distributed under the Creative Commons Attribution 4.0 International licence, as recorded in the arXiv licence metadata for this exact version and shown on the abstract page https://arxiv.org/abs/2605.25028v1. A full rights notice is delivered at licenses/arxiv-2605.25028-CC-BY-4.0.txt. Characters: every retained excerpt is pure ASCII, so the delivered engine, XeLaTeX with its default Latin Modern text font, reports no missing character.
\begin{lemma}\label{lem:p_poly}
  The volume of~$P(b)$, as a function in~$b\in\R^m$, is a polynomial
  of degree at most~$d$ on each box~$[\bar b,\bar b+e]$ with~$\bar
  b\in\Z^m$.
\end{lemma}

\begin{lemma}\label{lem:pseudo_vol}
  Let~$b\in\R^m$. If~$m$ is fixed, then the volume of~$P(b)$ can be
  computed in time polynomial in~$d$ and~$||A||_\infty$.
\end{lemma}

\begin{proof}
  For~$k\in\{0,1,\dots,d\}$ and~$s\in\R^m$,
  define~$P_{k}(s):=\{\xi\in[0,1]^k\mid A_k\xi\le s\}$,
  where~$A_k\in\Z^{m\times k}$ consists of the first~$k$ columns
  of~$A$. By Lemma~\ref{lem:p_poly}, the volume of~$P_{k}(s)$ is a
  polynomial in~$s$ of degree at most~$k$ on each box~$[\bar s,\bar
    s+e]$,~$\bar s\in\Z^m$. The at most~$k^m$ coefficients of these
  polynomials, for all~$\bar s$ in the set
  \[
  \textstyle Z_k:=\Z^m\cap [\lfloor
    b\rfloor-\sum_{\ell=k+1}^d||a_\ell||_\infty,\lfloor
    b\rfloor+\sum_{\ell=k+1}^d||a_\ell||_\infty]\;,
  \]  
  can be computed iteratively for~$k=0,\dots,d$ using dynamic
  programming: First, by convention,~$\vol(P_{0,s})=1$ if~$s\ge 0$ and
  zero otherwise, so that all polynomials for~$k=0$ can be computed
  easily. For~$k\ge 1$, we use the recursive formula
  \[
    \vol(P_{k}(s))=\int_0^1\vol(P_{k-1}(s-a_kt))\d t\quad\forall
    k\in\N,\,s\in\R^m\;.
  \]
  Again by Lemma~\ref{lem:p_poly}, the volume of~$P_{k-1}(s-a_kt)$ is
  a polynomial in~$t$ as long as~$s-a_kt$ remains in a box~$[\bar
    s,\bar s+e]$. To compute the coefficients of~$\vol(P_{k}(s))$
  on~$s\in[\bar s,\bar s+e]$ for some~$\bar s\in Z_k$, it suffices to
  symbolically integrate the polynomial~$\vol(P_{k-1}(s-a_kt))$, as a
  function in~$t$, for each of the corresponding ranges of~$t$, and to
  coefficient-wise sum up the resulting polynomials in~$s$ weighted by
  the lengths of these ranges. By definition of~$Z_k$, we have~$\bar
  s-a_kt\in[\bar{\bar s},\bar{\bar s}+e]$ for some~$\bar{\bar s}\in
  Z_{k-1}$, for all~$t\in[0,1]$, so that we know the coefficients
  of~$\vol(P_{k-1}(s-a_kt))$ from the previous iteration.

  Since~$|Z_k|\le (2(d-k)||A||_\infty+1)^m$ for all~$k$, all this can
  be done in time polynomial in~$d$ and~$||A||_\infty$ if~$m$ is
  fixed. Finally, the desired volume is obtained by evaluating the
  polynomial describing~$\vol(P_{d}(b))$ on~$[\lfloor b\rfloor,\lfloor
    b\rfloor+e]$ in the given~$b\in\R^m$.
\end{proof}

\end{document}
